\documentclass[12pt]{article}
\usepackage[utf8]{inputenc}
\usepackage[T1]{fontenc}
\usepackage[margin=1in]{geometry}
\usepackage{setspace}
\usepackage{xcolor}
\usepackage[colorlinks=true, urlcolor=blue!60!black, linkcolor=blue!60!black]{hyperref}
\setlength{\parindent}{2em}
\setlength{\parskip}{1.25em}
\renewcommand{\baselinestretch}{1.0}
\begin{document}
\begin{center}
{\LARGE \textit{Mathews v. Eldridge}, 424 U.S. 319 (1976)\par}
\vspace{0.8em}
{\large Case Brief --- Four-Page Edition\par}
\end{center}
\vspace{1em}

\section{Memory Jogger}

The Social Security Administration may cut off disability benefits before an evidentiary hearing. Due process is flexible, and courts weigh the private interest, the risk of error and value of added safeguards, and the government's interest; here those factors favored written pre-termination procedures plus a full post-termination hearing (\href{https://supreme.justia.com/cases/federal/us/424/319/}{424 U.S. 319}).

\section{Detailed Case Facts}

Social Security Disability Insurance, financed by payroll taxes, pays benefits to insured workers unable to engage in substantial gainful activity; unlike welfare, it is not means-tested. A state agency team monitors continuing eligibility through questionnaires and physicians' reports. If it tentatively finds the disability ended, it notifies the beneficiary in writing, summarizes the evidence, opens the file, and invites written rebuttal. After termination the claimant may obtain \textit{de novo} reconsideration, a full evidentiary hearing before an Administrative Law Judge, Appeals Council review, and judicial review, with full retroactive benefits if he prevails.

George Eldridge began receiving benefits in 1968 for chronic anxiety and back strain and was later diagnosed with diabetes. In 1972, after reviewing his questionnaire and reports from his physician and a psychiatric consultant, the Virginia agency tentatively found that his disability had ceased. Eldridge disputed the agency's description of his back condition in writing, but benefits were terminated. According to the dissent, the cutoff led to foreclosure on his home and repossession of his furniture.

\section{Procedural History}

Without seeking reconsideration, Eldridge sued, relying on \textit{Goldberg v.\ Kelly}, which required an oral evidentiary hearing before termination of welfare. The Western District of Virginia held his interest indistinguishable from a welfare recipient's and enjoined termination without a prior hearing; the Fourth Circuit affirmed. The Supreme Court reversed on February 24, 1976.

\section{Judicial Votes}

Reversed, \textbf{6--2}.
\begin{itemize}
\item \textbf{Majority:} Justice Powell, joined by Chief Justice Burger and Justices Stewart, White, Blackmun, and Rehnquist.
\item \textbf{Dissenting:} Justice Brennan, joined by Justice Marshall.
\item \textbf{Not participating:} Justice Stevens.
\end{itemize}

\section{Holding}

(1) The district court had jurisdiction under \S~405(g): presentment of a claim is nonwaivable, but exhaustion may be waived for a collateral constitutional claim with a colorable showing of irreparable harm. (2) Due process does not require an evidentiary hearing before termination of disability benefits. Courts balance (1) the private interest, (2) the risk of erroneous deprivation and the probable value of additional safeguards, and (3) the government's interest, including fiscal and administrative burdens. The SSA's written pre-termination procedures, combined with post-termination review and retroactive relief, satisfy the Fifth Amendment.

\section{Analysis of Opinions}

\textbf{Majority (Powell).} The Secretary conceded a property interest. Some hearing is required before a final deprivation, but due process is flexible and calls for such protections as the situation demands. \emph{Private interest:} because back pay is available, the interest is uninterrupted income pending decision. Unlike \textit{Goldberg}'s welfare recipients at the margin of subsistence, disability recipients may have other resources. Powell conceded that the distinction can be overstated, given delays of more than a year, but noted that other safety-net programs remain available.

\emph{Risk of error:} citing Judge Friendly, Powell explained that disability turns on routine, standard, and unbiased medical reports rather than the credibility questions in \textit{Goldberg}. Credibility may matter occasionally, but rules are shaped by the generality of cases. The questionnaire, open file, tentative determination, and invitation to rebut let claimants respond effectively. A 58.6\% reversal rate among appealed cases overstated error, since the overall error rate was far lower and new evidence was often introduced later.

\emph{Government's interest:} added hearings and payments to ineligible recipients impose costs, which may come out of the pockets of the deserving. Quoting Frankfurter, Powell warned against wholesale transplantation of judicial procedure into agencies and gave substantial weight to the judgments of Congress and administrators.

\textbf{Brennan (dissenting).} The author of \textit{Goldberg} would have required a pre-termination hearing. Calling the deprivation limited was speculative, as Eldridge's family sleeping in one bed showed; Congress's choice not to means-test implies need; and pointing to welfare is no argument for diluting due process.

\textbf{Later developments.} \textit{Loudermill}, \textit{Hamdi}, \textit{Turner v.\ Rogers}, and \textit{Nelson v.\ Colorado} applied \textit{Mathews}, while \textit{Medina v.\ California} and \textit{Culley v.\ Marshall} relied on history instead.

\section{Examples: Future Applications}

\begin{enumerate}
\item \textbf{Same-Side (Post-Hearing Suffices): A Veteran's Rating.} An agency reduces a veteran's non-means-tested disability rating on medical evidence after written notice and a chance to respond, with a later hearing and retroactive relief.
\item \textbf{Same-Side (No Pre-Suspension Hearing): A Charged Officer.} A university suspends a campus officer charged with a felony; the independent charge lowers the risk of error, as \textit{Gilbert v.\ Homar} held.
\item \textbf{Opposite-Side (Hearing Required): Welfare on a Caseworker's Suspicion.} Means-tested aid is cut off on a neighbor's report that turns on credibility; \textit{Goldberg} requires an oral pre-termination hearing.
\item \textbf{Opposite-Side (Opportunity Required): A Tenured Guard.} Under \textit{Loudermill}, a tenured employee is entitled to notice and an opportunity to respond before firing.
\item \textbf{Fence-Sitter: Algorithmic Medicaid Cuts.} A model proposes cutting in-home nursing hours, with individualized explanations and written responses. The State cites these safeguards; recipients reply that disputes over daily needs turn on credibility.
\end{enumerate}

\section{Critique}

\textbf{Dignitary values.} In \textit{Three Factors in Search of a Theory of Value}, Jerry Mashaw criticized \textit{Mathews} for reducing due process to utilitarian cost-benefit accounting while ignoring dignitary values, namely the intrinsic worth of in-person participation and reasoned explanation.

\textbf{History.} Originalists likewise fault \textit{Mathews} as an ahistorical balancing test, a view reflected in the Court's turn to historical practice in \textit{Medina} and \textit{Culley}.

\section{Key Quotations}

\begin{quote}
``More precisely, our prior decisions indicate that identification of the specific dictates of due process generally requires consideration of three distinct factors: first, the private interest that will be affected by the official action; second, the risk of an erroneous deprivation of such interest through the procedures used, and the probable value, if any, of additional or substitute procedural safeguards; and, finally, the Government's interest, including the function involved and the fiscal and administrative burdens that the additional or substitute procedural requirement would entail.''
\end{quote}
\noindent --- Justice Powell, Opinion of the Court, 424 U.S. at 334--35.

\begin{quote}
``To be sure, credibility and veracity may be a factor in the ultimate disability assessment in some cases. But procedural due process rules are shaped by the risk of error inherent in the truthfinding process as applied to the generality of cases, not the rare exceptions.''
\end{quote}
\noindent --- Justice Powell, Opinion of the Court, 424 U.S. at 344.

\begin{quote}
``The judicial model of an evidentiary hearing is neither a required, nor even the most effective, method of decisionmaking in all circumstances.''
\end{quote}
\noindent --- Justice Powell, Opinion of the Court, 424 U.S. at 348.

\end{document}
